\documentclass[12pt]{article}
\usepackage[utf8]{inputenc}
\usepackage[T1]{fontenc}
\usepackage[margin=1in]{geometry}
\usepackage{setspace}
\usepackage{xcolor}
\usepackage[colorlinks=true, urlcolor=blue!60!black, linkcolor=blue!60!black]{hyperref}
\setlength{\parindent}{2em}
\setlength{\parskip}{1.25em}
\renewcommand{\baselinestretch}{1.0}
\begin{document}
\begin{center}
{\LARGE \textit{Ex parte Young}, 209 U.S. 123 (1908)\par}
\vspace{0.8em}
{\large Case Brief --- Four-Page Edition\par}
\end{center}
\vspace{1em}

\section{Memory Jogger}

The Eleventh Amendment does not bar a federal court from enjoining a state official from enforcing an unconstitutional state statute: an officer who acts in conflict with the Federal Constitution is ``stripped of his official or representative character'' and treated as an individual rather than the State (\href{https://www.courtlistener.com/opinion/96819/ex-parte-young/}{CourtListener}).

\section{Detailed Case Facts}

During the Progressive Era, Minnesota sharply cut intrastate railroad rates and backed the cuts with severe penalties. A 1906 commission order lowered merchandise freight rates, with disobedience a gross misdemeanor carrying fines of up to \$10,000 per offense; an April 1907 act capped passenger fares at two cents per mile, making each overcharge a felony punishable by up to five years in prison. The railroads complied with both. Then Chapter 232, the Commodity Rate Act (effective June 1, 1907), set maximum carload freight rates and made any railroad officer or employee who violated it guilty of a misdemeanor punishable by up to 90 days in jail for each offense.

On May 31, 1907, shareholders of the Northern Pacific Railway sued the railway, the commissioners, and Attorney General Edward T. Young in federal circuit court, alleging that the cumulative rates were confiscatory under the Due Process Clause and that the draconian penalties effectively closed the courts to any test of the statutes' validity. They sought to bar the railway from adopting the rates and Young from enforcing them.

\section{Procedural History}

The circuit court issued a temporary restraining order, rejected Young's Eleventh Amendment objection, and on September 23, 1907, entered a preliminary injunction finding the commodity rates prima facie confiscatory. The next day Young defied it by obtaining a writ of mandamus from a Minnesota state court compelling the railway to adopt the rates. The federal court held him in contempt, fined him \$100, and committed him to the U.S. Marshal's custody. Young sought habeas corpus and certiorari in the Supreme Court, which decided the case on March 23, 1908.

\section{Judicial Votes}

Petition dismissed, \textbf{8--1}.
\begin{itemize}
\item \textbf{Majority:} Justice Peckham, joined by Chief Justice Fuller and Justices Brewer, White, McKenna, Holmes, Day, and Moody.
\item \textbf{Concurrences:} None.
\item \textbf{Dissent:} Justice Harlan.
\end{itemize}

\section{Holding}

The circuit court had jurisdiction to enjoin Young and to hold him in contempt. A federal court of equity may grant prospective injunctive relief against a state officer who has some duty connected with enforcing an allegedly unconstitutional statute, whether that duty comes from the statute itself or from general state law. Because the State cannot authorize its officers to violate the Constitution, the officer cannot claim the State's immunity. Penalties so severe as to deter parties from testing a rate law in court are unconstitutional on their face. Peckham's statements that federal courts may not enjoin state courts or grand juries are influential dicta.

\section{Analysis of Opinions}

\textbf{Majority (Peckham).} Confiscatory rates and penalties that deter access to the courts raise substantial federal questions. Subjecting employees to imprisonment and ruinous fines merely for testing a technically complex rate schedule closed off all approaches to the courts. On the Eleventh Amendment, Peckham reaffirmed \textit{Hans v. Louisiana} but, drawing on \textit{Osborn} through \textit{Smyth v. Ames}, held that an officer enforcing an unconstitutional law acts without valid state authority. He distinguished \textit{Fitts v. McGhee}, where the officers had no enforcement connection: Young's general duty to enforce Minnesota law sufficed. An officer has no lawful discretion to enforce an unconstitutional act, and awaiting prosecution was an inadequate remedy given the crushing penalties.

\textbf{Harlan (dissenting).} The distinction between the State and its Attorney General is a fiction that would ``practically obliterate the Eleventh Amendment.'' A State can act only through its officers, so forbidding its chief law officer from invoking its own courts is a suit against the State. The ruling departed from \textit{In re Ayers}, ``frittered away'' \textit{Fitts}, and would reduce States to ``dependencies.'' State courts share the oath to enforce federal rights, and the railway should have raised its defenses there, subject to Supreme Court review.

\section{Examples: Future Applications}

\begin{enumerate}
\item \textbf{Same-Side (Permitted): Clinic Licensing.} A clinic sues the health commissioner charged with revoking licenses and imposing \$50,000 daily fines under a law banning a federally protected procedure. A direct enforcement connection and prospective relief fit \textit{Young}.
\item \textbf{Same-Side (Permitted): Clean Air Act Compliance.} A utility obtains an injunction requiring a state agency director to follow federal permitting standards, even though compliance costs money; under \textit{Edelman v. Jordan}, an ``ancillary effect'' on the treasury is permissible.
\item \textbf{Opposite-Side (Barred): Retroactive Benefits.} An order compelling payment of two years of withheld welfare benefits from the state treasury is indistinguishable from damages, as \textit{Edelman} held.
\item \textbf{Opposite-Side (Barred): State-Law Claims.} An injunction requiring officials to obey a \emph{state} mental-health statute is barred under \textit{Pennhurst}, because \textit{Young} vindicates the supremacy of federal law.
\item \textbf{Fence-Sitter: Private Enforcement.} A state forbids its officials to enforce a law restricting a constitutional right and delegates enforcement to private suits. As in \textit{Whole Woman's Health v. Jackson}, suits against judges, clerks, and the attorney general fail, while claims against licensing officials may survive.
\end{enumerate}

\section{Critique}

\textbf{The double fiction.} Young's enforcement is ``state action'' for the Fourteenth Amendment yet an unauthorized private wrong for the Eleventh. Formalist critics echo Harlan, and Judge Andrew Oldham and coauthors argue that 1789 English Chancery would not have recognized \textit{Young}'s equitable cause of action. John Harrison, by contrast, defends \textit{Young} as a traditional anti-suit injunction, and James Pfander and Jacob Wentzel trace officer suits to deep equitable traditions.

\textbf{A reversal of reputation.} Progressives first denounced \textit{Young} as ``government by injunction,'' prompting Congress in 1910 to require three-judge courts; as Edward Purcell observes, it later became the backbone of civil rights litigation. Because it rests on officer-suit formalism, the Court has narrowed it in \textit{Edelman}, \textit{Seminole Tribe}, and \textit{Whole Woman's Health}.

\section{Key Quotations}

\begin{quote}
``If the act which the state Attorney General seeks to enforce be a violation of the Federal Constitution, the officer, in proceeding under such enactment, comes into conflict with the superior authority of that Constitution, and he is, in that case, stripped of his official or representative character, and is subjected in his person to the consequences of his individual conduct. The State has no power to impart to him any immunity from responsibility to the supreme authority of the United States.''
\end{quote}
\noindent --- Justice Peckham, Opinion of the Court, 209 U.S. at 159--60.

\begin{quote}
``In making an officer of the State a party defendant in a suit to enjoin the enforcement of an act alleged to be unconstitutional, it is plain that such officer must have some connection with the enforcement of the act, or else it is merely making him a party as a representative of the State, and thereby attempting to make the State a party.''
\end{quote}
\noindent --- Justice Peckham, Opinion of the Court, 209 U.S. at 157.

\begin{quote}
``We hold, therefore, that the provisions of the acts relating to the enforcement of the rates, either for freight or passengers, by imposing such enormous fines and possible imprisonment as a result of an unsuccessful effort to test the validity of the laws themselves, are unconstitutional on their face, without regard to the question of the insufficiency of those rates.''
\end{quote}
\noindent --- Justice Peckham, Opinion of the Court, 209 U.S. at 148.

\end{document}
